\documentclass[6.5pt,landscape,a4paper]{extarticle}
\usepackage[landscape,a4paper,top=2mm,bottom=2mm,left=2mm,right=2mm]{geometry}
\usepackage{multicol}
\usepackage{amsmath,amssymb}
\usepackage[T1]{fontenc}


\usepackage{xcolor}
\usepackage{booktabs}
\usepackage{array}
\usepackage{enumitem}
\usepackage{mdframed}
\usepackage{tabularx}
\usepackage[colorlinks=true,urlcolor=red]{hyperref}


% ---- layout ----
\setlength{\columnsep}{2mm}
\setlength{\columnseprule}{0.25pt}
\setlength{\parindent}{0pt}
\setlength{\parskip}{0.4pt}
\setlength{\abovedisplayskip}{1pt}
\setlength{\belowdisplayskip}{1pt}
\setlength{\abovedisplayshortskip}{0pt}
\setlength{\belowdisplayshortskip}{0pt}
\pagestyle{empty}

% ---- colours ----
\definecolor{hdr}{RGB}{20,60,140}
\definecolor{sub}{RGB}{60,100,180}
\definecolor{keyc}{RGB}{230,240,255}
\definecolor{warnc}{RGB}{255,235,235}

% ---- section headers ----
\makeatletter
\renewcommand{\section}{\@startsection{section}{1}{0mm}%
  {-1pt plus -0.3pt}{0.3pt plus 0.1pt}%
  {\normalfont\bfseries\small\color{white}%
   \setlength{\fboxsep}{1pt}\colorbox{hdr}}}
\renewcommand{\subsection}{\@startsection{subsection}{2}{0mm}%
  {-0.8pt plus -0.3pt}{0.2pt plus 0.1pt}%
  {\normalfont\bfseries\footnotesize\color{sub}}}
\makeatother
\setcounter{secnumdepth}{0}

% ---- compact lists ----
\setlist{nosep,leftmargin=*,topsep=0pt,itemsep=0pt,parsep=0pt}

% ---- key result box ----
\newmdenv[topline=false,bottomline=false,rightline=false,
  linewidth=1.8pt,linecolor=hdr,
  innertopmargin=1pt,innerbottommargin=1pt,
  innerleftmargin=3pt,innerrightmargin=1pt,
  skipabove=1pt,skipbelow=1pt,
  backgroundcolor=keyc]{kbox}

% ---- warn box ----
\newmdenv[topline=false,bottomline=false,rightline=false,
  linewidth=1.8pt,linecolor=red!70!black,
  innertopmargin=1pt,innerbottommargin=1pt,
  innerleftmargin=3pt,innerrightmargin=1pt,
  skipabove=1pt,skipbelow=1pt,
  backgroundcolor=warnc]{wbox}



\begin{document}

% ---- title bar ----
\noindent\colorbox{hdr}{\parbox{\dimexpr\linewidth-2\fboxsep}{%
    \color{white}\bfseries\small
    QF1100 Introduction to Quantitative Finance --- Final Exam Cheatsheet \href{https://github.com/stinkray77}{@stinkray77}
    \hfill \quad AY2025/2026 Sem 2
  }}
\vspace{0.1mm}

\begin{multicols*}{4}
  \scriptsize

  %% =====================================================================
  \section{1.\ Theory of Interest}

  \subsection{Accumulation Functions}
  \textbf{Simple:} \(a(t)=1+rt\)\quad\textbf{Compound:} \(a(t)=(1+r)^t\) \textsf{[default]}

  \textbf{Continuous:} \(a(t)=e^{\delta t}\)

  \textbf{Accumulated value} of principal \(P\): \(A(t)=P\cdot a(t)\)

  \subsection{Frequency of Compounding}
  Nominal rate \(r^{(p)}\) paid \(p\) times/yr:
  \[a(t)=\Bigl(1+\tfrac{r^{(p)}}{p}\Bigr)^{pt}\]
  Effective annual rate:
  \[r_e=\Bigl(1+\tfrac{r^{(p)}}{p}\Bigr)^p - 1\]
  \textbf{Equivalence} of two nominal rates:
  \[\Bigl(1+\tfrac{r^{(p)}}{p}\Bigr)^p = \Bigl(1+\tfrac{R^{(q)}}{q}\Bigr)^q\]

  \begin{kbox}
    \(r_e > r^{(p)}\) for any \(p>1\). Continuous compounding maximises \(r_e\) for a given nominal rate.
  \end{kbox}

  \subsection{Force of Interest}
  \[\delta(t)=\frac{a'(t)}{a(t)}=\frac{d}{dt}\ln a(t)\]
  Constant $\delta$: $a(t)=e^{\delta t}$, $\delta=\ln(1+r_e)$, $r_e=e^\delta-1$.

  \begin{kbox}
    \textbf{Variable force of interest:}
    \[a(t)=\exp\!\Bigl(\int_0^t\!\delta(u)\,du\Bigr)\quad a(s,t)=\exp\!\Bigl(\int_s^t\!\delta(u)\,du\Bigr)\]
  \end{kbox}

  \subsection{Key Conversions (quick ref)}
  \begin{tabular}{@{}ll@{}}
    \toprule
    Convert                & Formula                                    \\
    \midrule
    \(r^{(p)}\to r_e\)     & \((1+r^{(p)}/p)^p-1\)                      \\
    \(r_e\to r^{(p)}\)     & \(p[(1+r_e)^{1/p}-1]\)                     \\
    \(r_e\to\delta\)       & \(\ln(1+r_e)\)                             \\
    \(\delta\to r_e\)      & \(e^{\delta}-1\)                           \\
    \(r^{(p)}\to r^{(q)}\) & equate \((1+r^{(p)}/p)^p=(1+r^{(q)}/q)^q\) \\
    \bottomrule
  \end{tabular}

  %% =====================================================================
  \section{2.\ Present Value \& Cash Flows}

  \subsection{Present Value (single payment)}
  \[\mathrm{PV}=\frac{c}{a(t)}=\frac{c}{(1+r)^t}\]
  Discount factor: \(v=\frac{1}{1+r}\), so \(\mathrm{PV}=c\cdot v^t\)

  \subsection{PV of Cash Flow Stream}
  \(C=\{(c_i,t_i)\}_{i=1}^n\):
  \[\mathrm{PV}(C)=\sum_{i=1}^n\frac{c_i}{(1+r)^{t_i}}\]
  \textbf{Time value} at \(t\): \(\mathrm{TV}_t(C)=\mathrm{PV}(C)\cdot(1+r)^t\)

  \subsection{Annuities}
  Let \(r\) = effective rate per period, \(A\) = payment per period, \(n\) = number of periods.

  \textbf{Ordinary annuity} (payment at \emph{end} of each period):
  \[\mathrm{PV}=\frac{A}{r}\Bigl[1-\frac{1}{(1+r)^n}\Bigr]\]

  \textbf{Annuity-due} (payment at \emph{start} of each period):
  \[\mathrm{PV}=\frac{A(1+r)}{r}\Bigl[1-\frac{1}{(1+r)^n}\Bigr]\]
  \(\Rightarrow\) PV(due) \(=\) PV(ordinary) \(\times(1+r)\)

  \textbf{Perpetuity} (ordinary, payments at end):
  \[\mathrm{PV}=\frac{A}{r}\]

  \textbf{Perpetuity-due} (payments at start):
  \[\mathrm{PV}=\frac{A(1+r)}{r}\]

  \textbf{Future value} (ordinary annuity at \(t=n\)):
  \[\mathrm{FV}=\frac{A}{r}\bigl[(1+r)^n-1\bigr]\]

  \textbf{Deferred annuity} ($n$ payments, first at $t=m+1$):
  \[\mathrm{PV}=\frac{A}{r}\Bigl[1-\frac{1}{(1+r)^n}\Bigr]\cdot\frac{1}{(1+r)^m}\]
  Discount ordinary annuity PV back by $m$ periods.

  \begin{wbox}
    \textbf{Rate mismatch:} Always convert nominal rate to \emph{effective rate per payment period} before applying annuity formula. E.g.\ 6\% p.a.\ convertible semi-annually \(\to\) effective monthly rate \(=(1+0.03)^{1/6}-1\).
  \end{wbox}

  \subsection{NPV \& IRR Decision Rules}
  \textbf{NPV criterion:} Invest if $\mathrm{NPV}>0$. Prefer A over B if $\mathrm{NPV}(A)>\mathrm{NPV}(B)$.

  \textbf{IRR criterion:} Invest if $\mathrm{IRR}>$ required return. Prefer A over B if $\mathrm{IRR}(A)>\mathrm{IRR}(B)$.

  IRR is $r$ s.t.\ $\mathrm{PV}(C)=0$. \textbf{Uniqueness:} if $c_0<0$, $c_k\ge0\;\forall k\ge1$, $\sum c_k>0 \Rightarrow$ unique IRR.

  Newton--Raphson: $r_{n+1}=r_n - f(r_n)/f'(r_n)$

  \subsection{Growing Annuity \& Perpetuity}
  First payment $A$ at $t=1$, grows at rate $g$/period, $n$ periods:
  \[\mathrm{PV}=\frac{A}{r-g}\Bigl[1-\Bigl(\frac{1+g}{1+r}\Bigr)^n\Bigr]\quad(r\neq g)\qquad r=g:\frac{nA}{1+r}\]
  Growing perpetuity ($r>g$): $\mathrm{PV}=A/(r-g)$.

  \subsection{Loans, Refinancing \& Last Payment}
  Loan $L$, $n$ payments $A$, rate $r$/period:
  $L=\frac{A}{r}[1-(1+r)^{-n}]$,\quad $A=\frac{Lr}{1-(1+r)^{-n}}$

  \textbf{Balance after $m$ payments:}
  $L_m^{\text{bal}}=\frac{A}{r}[1-(1+r)^{-(n-m)}]$ (prospective)

  $L_m^{\text{bal}}=L(1+r)^m-\frac{A}{r}[(1+r)^m-1]$ (retrospective)

  \textbf{Irregular last payment} ($n$ non-integer): make $\lfloor n\rfloor$ full payments of $A$ plus final smaller payment $B$.
  \begin{itemize}
    \item $B$ at $t=\lfloor n\rfloor$: set $L=\mathrm{PV}$ of $\lfloor n\rfloor{-}1$ payments of $A$ plus $B$ at $t=\lfloor n\rfloor$.
    \item $B$ at $t=\lfloor n\rfloor+1$: use $L_{\lfloor n\rfloor}^{\text{bal}}$, then $B=L_{\lfloor n\rfloor}^{\text{bal}}(1+r)$.
  \end{itemize}

  %% =====================================================================
  \section{3.\ Bonds}

  \subsection{Notation}
  \begin{tabular}{@{}ll@{}}
    \(F\)       & face value (= \(R\) unless stated) \\
    \(R\)       & redemption value                   \\
    \(c\)       & nominal coupon rate (annual)       \\
    \(m\)       & coupon payments per year           \\
    \(n\)       & total coupon payments              \\
    \(\lambda\) & nominal bond yield (annual)        \\
    \(P\)       & bond price                         \\
  \end{tabular}

  \subsection{Bond Pricing Formula}
  \begin{kbox}
    \[P=\frac{R}{(1+\lambda/m)^n}+\sum_{i=1}^n\frac{cF/m}{(1+\lambda/m)^i}\]
    \[=F\Biggl[\frac{1}{(1+\lambda/m)^n}+\frac{c/m}{\lambda/m}\Bigl(1-\frac{1}{(1+\lambda/m)^n}\Bigr)\Biggr]\quad(F=R)\]
  \end{kbox}

  \textbf{Premium/discount form} ($F=R$):
  $P=F+F\cdot\frac{c-\lambda}{\lambda}\bigl[1-(1+\lambda/m)^{-n}\bigr]$

  $c>\lambda\Rightarrow P>F$ (premium);\quad $c=\lambda\Rightarrow P=F$ (par);\quad $c<\lambda\Rightarrow P<F$ (discount).

  \begin{kbox}
    \textbf{Makeham formula} (alternative; let $K=R/(1+\lambda/m)^n$):
    \[P=K+\frac{c}{\lambda}(F-K)\qquad(F=R\text{ case})\]
  \end{kbox}

  \textbf{Recursive price:} $P_{k+1}=P_k(1+\lambda/m)-cF/m$ (price after $k$-th coupon).

  \subsection{Variable Coupon Bonds}
  Split into groups with constant coupon rate; price each group as a separate annuity + shared redemption PV.

  \begin{wbox}
    \textbf{Exam pattern (Q3 midterm):} Bond with different coupon rates for each 5-year block and \(R\neq F\). Price each block as an annuity discounted at yield \(\lambda\), then add PV of redemption at maturity.
  \end{wbox}

  \subsection{Determining Bond Yield (Newton--Raphson)}
  $f(\lambda)=P(\lambda)-P_{\text{market}}$, iterate $\lambda_{n+1}=\lambda_n - f(\lambda_n)/f'(\lambda_n)$.

  \[P(\lambda)=\frac{R}{(1+\lambda/m)^n}+\frac{cF/m}{\lambda/m}\Bigl[1-\frac{1}{(1+\lambda/m)^n}\Bigr]\]

  \begin{kbox}
    $f'(\lambda)$ for annual coupon bond ($m=1$, $F=R$):
    \[f'(\lambda)=-\sum_{i=1}^{n}\frac{i\cdot cF}{(1+\lambda)^{i+1}}-\frac{nF}{(1+\lambda)^{n+1}}=-\frac{D\cdot P(\lambda)}{1+\lambda}\]
    where $D$ = Macaulay duration. (In general: $f'(\lambda)=-\frac{D\cdot P}{m(1+\lambda/m)}$.)
  \end{kbox}

  Start: $\lambda_0=c$. Price--yield: \textbf{inverse} \& \textbf{convex} ($P\downarrow$ as $\lambda\uparrow$).

  \subsection{Macaulay Duration}
  \[D=\frac{\sum_{i=1}^n t_i\cdot\mathrm{PV}(c_i,t_i)}{\mathrm{PV}(C)}=\sum_{i=1}^n w_i t_i,\quad w_i=\frac{\mathrm{PV}(c_i)}{\mathrm{PV}(C)}\]

  \textbf{Special cases} ($\mu=\lambda/m$, $\gamma=c/m$):
  Zero-coupon: $D=T$.\quad Perpetuity : $D=(1+\lambda / m)/\lambda$.

  \begin{kbox}
    \textbf{Coupon bond, $F=R$ (general):}
    $D=\dfrac{1+\mu}{m\mu}-\dfrac{1+\mu+n(\gamma-\mu)}{m\gamma[(1+\mu)^n-1]+m\mu}$\quad\textbf{(2)}

    \textbf{At par $P=F=R$ ($\mu=\gamma$):}
    $D=\dfrac{1+\mu}{m\mu}\Bigl[1-\dfrac{1}{(1+\mu)^n}\Bigr]$\quad\textbf{(3)}
  \end{kbox}

  $t_1\le D\le t_n$;\quad $D$ independent of $F$ and of yield compounding frequency.

  %% =====================================================================
  \section{4.\ Term Structure}

  \subsection{Spot Rates \& Forward Rates}
  \textbf{Spot rate} \(s_t\): effective annual rate from 0 to \(t\).
  Accumulation: \((1+s_t)^t\).

  \textbf{Forward rate} \(f_{j,k}\): rate from \(t=j\) to \(t=k\) (annual, same convention as spot).

  \begin{kbox}
    \textbf{No-arbitrage relation:}
    \[(1+s_k)^k=(1+s_j)^j(1+f_{j,k})^{k-j}\]
    Note: \(f_{0,t}=s_t\).
  \end{kbox}

  \textbf{Continuous compounding version:}
  \[e^{s_k k}=e^{s_j j}\cdot e^{f_{j,k}(k-j)}\quad\Rightarrow\quad f_{j,k}=\frac{s_k k - s_j j}{k-j}\]

  \subsection{Bootstrapping Spot Rates}
  Use observed coupon bond prices to extract \(s_1,s_2,\ldots\) sequentially:

  For a coupon bond with annual payments:
  \[P=\frac{c_1}{1+s_1}+\frac{c_2}{(1+s_2)^2}+\cdots+\frac{c_n+F}{(1+s_n)^n}\]
  Solve for the unknown \(s_n\) given \(s_1,\ldots,s_{n-1}\) already found.

  \begin{wbox}
    \textbf{Exam pattern:} Extract \(s_1,s_2\) from zero-coupon bond prices, then use coupon bond price to find \(s_3\), then compute \(f_{1,3}\).
  \end{wbox}

  \subsection{Yield Curve Shapes}
  Normal (upward): economic expansion, low inflation.\\
  Inverted (downward): pre-recession, high inflation.\\
  Flat term structure: all $s_t=r$ equal. \textbf{Proof} ($\delta$ constant): continuous compounding $\Rightarrow a(t)=e^{rt}$, so $\delta(t)=a'(t)/a(t)=re^{rt}/e^{rt}=r$. Semiannual: $(1+s_t/2)^2=e^{r_t}$; if $s_t$ constant then $r_t$ constant and vice versa. $\checkmark$

  %% =====================================================================
  \section{5.\ Asset Returns \& Portfolios}

  \subsection{Returns}
  \textbf{Total return:} \(R=P_1/P_0\)\quad
  \textbf{Rate of return:} \(r=(P_1-P_0)/P_0\)\quad\(R=1+r\)

  \textbf{Short sell:} \(R_{\text{short}}=P_1/P_0\) (same formula -- negatives cancel).

  \subsection{Definitions \& Identities}
  $\text{Var}(X)=E[X^2]-E[X]^2$\quad $\text{Cov}(X,Y)=E[XY]-E[X]E[Y]$\quad $\rho_{XY}=\dfrac{\text{Cov}(X,Y)}{\sigma_X\sigma_Y}$

  \textbf{Expectation (linearity):} $E(c)=c$;\; $E(cX)=cE(X)$;\; $E(X+c)=E(X)+c$;\; $E(aX+bY)=aE(X)+bE(Y)$

  \textbf{Variance:} $\text{Var}(aX+b)=a^2\text{Var}(X)$;\quad $\text{Var}(X\pm Y)=\text{Var}(X)+\text{Var}(Y)\pm2\text{Cov}(X,Y)$

  \textbf{Covariance:} $\text{Cov}(X,Y)=\text{Cov}(Y,X)$;\; $\text{Cov}(X,X)=\text{Var}(X)$;\; $\text{Cov}(aX,bY)=ab\,\text{Cov}(X,Y)$

  \textbf{Bilinearity:} $\text{Cov}(aX+bY,Z)=a\,\text{Cov}(X,Z)+b\,\text{Cov}(Y,Z)$

  Indep $\Rightarrow$ $\text{Cov}=0$;\quad $\text{Cov}=0 \not\Rightarrow$ indep.

  \subsection{Asset Returns (notation)}
  Mean: $\mu_i=E(r_i)$;\; Var: $\sigma_i^2$;\; Cov: $\sigma_{ij}=\text{Cov}(r_i,r_j)$\quad\textbf{Short-sell}: profit $=P_0-P_1$, return $=-r_{\text{long}}$, loss unlimited.

  \subsection{Portfolio of 2 Assets}
  Weights \(w_1=\alpha\), \(w_2=1-\alpha\) (\(\alpha\in\mathbb{R}\) if short-selling allowed; \(\alpha\in[0,1]\) otherwise).

  \[\mu_p=\alpha\mu_1+(1-\alpha)\mu_2\]
  \[\sigma_p^2=\alpha^2\sigma_1^2+(1-\alpha)^2\sigma_2^2+2\alpha(1-\alpha)\rho_{12}\sigma_1\sigma_2\]

  \subsection{Global Minimum-Variance Portfolio (GMVP)}
  \begin{kbox}
    \[\alpha^*=\frac{\sigma_2^2-\rho_{12}\sigma_1\sigma_2}{\sigma_1^2+\sigma_2^2-2\rho_{12}\sigma_1\sigma_2}\]
    \[(\sigma_p^2)^*=\frac{\sigma_1^2\sigma_2^2(1-\rho_{12}^2)}{\sigma_1^2+\sigma_2^2-2\rho_{12}\sigma_1\sigma_2}\]
  \end{kbox}

  If \(\rho=1\): feasible set is a straight line. If \(\rho=-1\): V-shape (zero variance achievable). If \(-1<\rho<1\): curved bullet shape.

  \textbf{Efficient frontier}: upper half of feasible set (including GMVP).

  \subsection{Constrained GMVP ($\mu_p\ge\bar\mu$)}
  If $\bar\mu>\mu_{\text{GMVP}}$, constraint is binding: set $\mu_p=\bar\mu$, solve $\alpha=\frac{\bar\mu-\mu_2}{\mu_1-\mu_2}$, substitute into $\sigma_p^2$.

  \textbf{No short-selling} ($\alpha\in[0,1]$): if unconstrained $\alpha^*\notin[0,1]$, GMVP is the asset with lower $\sigma$.

  \begin{wbox}
    Check: if \(\alpha^*\notin[0,1]\) and short-selling not allowed, the GMVP is the asset with lower \(\sigma\).
  \end{wbox}

  %% =====================================================================
  \section{6.\ Forwards \& Futures}

  \subsection{Forward Contract}
  Obligation to buy/sell at forward price \(F(t,T)\) at maturity \(T\).
  \begin{itemize}
    \item Long payoff: \(S_T-F(t,T)\)
    \item Short payoff: \(F(t,T)-S_T\)
    \item Costs nothing to enter (no premium)
  \end{itemize}

  \subsection{Forward Pricing (No-Arbitrage)}
  \begin{kbox}
    \textbf{Non-dividend-paying asset:}
    \[F(0,T)=S_0 e^{rT}\ (\text{cts})\quad F(0,T)=S_0(1+r)^T\ (\text{disc})\]

    \textbf{General \(t\):} \(F(t,T)=S_t\,e^{r(T-t)}\)

    \textbf{With storage cost} (lump sum $c$ at $t=0$):
    \[F(0,T)=(S_0+c)\,e^{rT}\]

    \textbf{With periodic storage} (payments $c$ at $t=0,\Delta,2\Delta,\ldots$): compute directly as annuity-due discounted to $t=0$.
  \end{kbox}

  \begin{wbox}
    Forward price does \textbf{not} depend on expected return \(\mu\) of the underlying --- only on \(S_0\), \(r\), \(T\) (and carry costs). This is a no-arbitrage result, not an expectation.
  \end{wbox}

  \subsection{Arbitrage Construction}
  If \(F(0,T)>S_0e^{rT}\): short forward, buy asset (borrow \(S_0\)). Profit at \(T\): \(F-S_0e^{rT}>0\).\\
  If \(F(0,T)<S_0e^{rT}\): long forward, short-sell asset (invest proceeds). Profit at \(T\): \(S_0e^{rT}-F>0\).

  \subsection{Futures vs.\ Forwards}
  \begin{tabular}{@{}lll@{}}
    \toprule
                 & Forward     & Futures   \\
    \midrule
    Settlement   & At maturity & Daily MTM \\
    Counterparty & Bilateral   & Exchange  \\
    \bottomrule
  \end{tabular}

  Daily P\&L long futures: \(\approx F_{t+1}-F_t\) (margin account).

  \subsection{Margin Account \& Mark-to-Market}
  \textbf{Setup:} $N$ contracts, contract size $Q$, futures price $F_t$.

  \textbf{Daily MTM:} Gain/loss $=(F_t-F_{t-1})\times Q\times N$ (positive for long, negative for short).

  \textbf{Margin call:} If balance $<$ maintenance margin $\Rightarrow$ top up to \emph{initial} margin. Top-up $=$ initial $-$ current balance.

  \textbf{Total contract value} $=F_0\times Q\times N$.\quad Initial margin $=$ fraction $\times$ TCV.\quad Maintenance margin $=$ fraction $\times$ TCV.

  \begin{wbox}
    Top up to \textbf{initial}, not maintenance. Missing this costs full marks.
  \end{wbox}

  %% =====================================================================
  \section{7.\ Hedging}

  \subsection{Perfect Hedge (Forwards/Futures)}
  Enter opposite position to fix price and eliminate risk. Conditions for perfect futures hedge: (1) delivery date matches; (2) same underlying; (3) integer multiple of contract size; (4) sufficient liquidity.

  \subsection{Hedging Direction ($\rho>0$)}
  {\tiny\begin{tabular}{@{}llll@{}}
      \toprule
      Obligation   & Risk               & Action              & P\&L           \\
      \midrule
      Buy asset A  & Price $\uparrow$   & Long Futures $A/B$  & $+(F_T{-}F_0)$ \\
      Sell asset A & Price $\downarrow$ & Short Futures $A/B$ & $+(F_0{-}F_T)$ \\
      \bottomrule
    \end{tabular}}\\
  If $\rho_{A,B}<0$: reverse direction.

  \subsection{Minimum-Variance Hedge}
  Hedge \(W\) units of asset A using futures on asset B (correlated). Short \(h\) futures contracts.

  Cash flow at \(T\): \(y=-WS_T+(F_T-F_0)h\)

  \[\text{Var}(y)=W^2\text{Var}(S_T)-2W\text{Cov}(S_T,F_T)h+\text{Var}(F_T)h^2\]

  \begin{kbox}
    \textbf{Optimal hedge ratio:}
    \[h^*=W\cdot\frac{\text{Cov}(S_T,F_T)}{\text{Var}(F_T)}=W\cdot\frac{\rho_{S,F}\,\sigma_S}{\sigma_F}\]

    \textbf{Minimum variance:}
    \[\text{Var}^*(y)=W^2\text{Var}(S_T)(1-\rho_{S,F}^2)\]

    Hedge effectiveness: \(\rho_{S,F}^2\). Factor \(\sqrt{1-\rho^2}\) measures residual risk.
  \end{kbox}

  \begin{wbox}
    \(h^*\) is in \emph{units}, not contracts. Number of contracts \(= h^*/(\text{contract size})\). Round to integer for real problems.
  \end{wbox}

  %% =====================================================================
  \section{8.\ Options}

  \subsection{Key Definitions}
  \textbf{Call}: right to \emph{buy} at strike \(K\).\quad\textbf{Put}: right to \emph{sell} at strike \(K\).

  Long = holder (pays premium); Short = writer (receives premium).

  European: exercise only at \(T\). American: exercise any time \(\le T\). (Assume European unless stated.)

  \subsection{Payoffs \& Profits}
  \begin{tabular}{@{}lll@{}}
    \toprule
    Position   & Payoff at \(T\)    & Profit               \\
    \midrule
    Long call  & \(\max(S_T-K,0)\)  & payoff \(-C e^{rT}\) \\
    Short call & \(-\max(S_T-K,0)\) & payoff \(+C e^{rT}\) \\
    Long put   & \(\max(K-S_T,0)\)  & payoff \(-P e^{rT}\) \\
    Short put  & \(-\max(K-S_T,0)\) & payoff \(+P e^{rT}\) \\
    \bottomrule
  \end{tabular}

  \(C, P\) = option premium (time-0 price). Time-value of premium at \(T\): \(Ce^{rT}\) or \(P e^{rT}\).

  \subsection{Bounds}
  \textbf{Upper bounds:} \(C\le S_0\)\qquad\(P\le Ke^{-rT}\)

  \textbf{Lower bounds:}
  \[C\ge\max(0,\;S_0-Ke^{-rT})\qquad P\ge\max(0,\;Ke^{-rT}-S_0)\]

  \subsection{Bounds Arbitrage}
  \textbf{$C>S_0$:} short call, long stock.\\
  {\tiny\begin{tabular}{@{}llll@{}}
    \toprule
                   & $t{=}0$   & $S_T{>}K$  & $S_T{\le}K$ \\
    \midrule
    Short call     & $+C$      & $-(S_T-K)$ & $0$         \\
    Long stock     & $-S_0$    & $+S_T$     & $+S_T$      \\
    \midrule
    \textbf{Total} & $C-S_0>0$ & $+K>0$     & $S_T\ge0$   \\
    \bottomrule
  \end{tabular}}

  \textbf{$P>Ke^{-rT}$:} short put, invest $Ke^{-rT}$.\\
  {\tiny\begin{tabular}{@{}llll@{}}
    \toprule
                   & $t{=}0$        & $S_T{<}K$  & $S_T{\ge}K$ \\
    \midrule
    Short put      & $+P$           & $-(K-S_T)$ & $0$         \\
    Bond           & $-Ke^{-rT}$    & $+K$       & $+K$        \\
    \midrule
    \textbf{Total} & $P-Ke^{-rT}>0$ & $S_T\ge0$  & $+K>0$      \\
    \bottomrule
  \end{tabular}}

  \textbf{$C<S_0-Ke^{-rT}$:} short stock, long call, invest $S_0-C$ [$I\equiv(S_0{-}C)e^{rT}$].\\
  {\tiny\begin{tabular}{@{}llll@{}}
    \toprule
                   & $t{=}0$      & $S_T{\ge}K$  & $S_T{<}K$ \\
    \midrule
    Short stock    & $+S_0$       & $-S_T$       & $-S_T$    \\
    Long call      & $-C$         & $+(S_T{-}K)$ & $0$       \\
    Invest         & $-(S_0{-}C)$ & $+I$         & $+I$      \\
    \midrule
    \textbf{Total} & $0$          & $I-K>0$      & $I-S_T>0$ \\
    \bottomrule
  \end{tabular}}

  \textbf{$P<Ke^{-rT}-S_0$:} long put, long stock, borrow $S_0+P$ [$B\equiv(S_0{+}P)e^{rT}$].\\
  {\tiny\begin{tabular}{@{}llll@{}}
    \toprule
                   & $t{=}0$      & $S_T{<}K$    & $S_T{\ge}K$ \\
    \midrule
    Long put       & $-P$         & $+(K{-}S_T)$ & $0$         \\
    Long stock     & $-S_0$       & $+S_T$       & $+S_T$      \\
    Borrow         & $+(S_0{+}P)$ & $-B$         & $-B$        \\
    \midrule
    \textbf{Total} & $0$          & $K-B>0$      & $S_T-B\ge0$ \\
    \bottomrule
  \end{tabular}}

  \begin{kbox}
    \textbf{Put--Call Parity} (European, no dividends):
    \[C + Ke^{-rT} = P + S_0\]
    {\tiny\begin{tabular}{@{}lllll@{}}
      \toprule
               &                & $t{=}T$   &           \\
               &                & $S_T{>}K$ & $S_T{<}K$ \\
      \midrule
      Port.\ A & Call           & $S_T-K$   & $0$       \\
               & ZC bond ($K$)  & $K$       & $K$       \\
               & \textit{Total} & $S_T$     & $K$       \\
      \midrule
      Port.\ C & Put            & $0$       & $K-S_T$   \\
               & Share          & $S_T$     & $S_T$     \\
               & \textit{Total} & $S_T$     & $K$       \\
      \bottomrule
    \end{tabular}}\\
    Same payoff $\Rightarrow$ $PV_A=PV_C$: $C+dK=P+S_0$.
  \end{kbox}

  \subsection{Put--Call Parity (General Form)}
  Let $d$ = discount factor. Discrete: $d=(1+r)^{-T}$. Continuous: $d=e^{-rT}$.
  \begin{kbox}
    \[C + dK = P + S_0\]
  \end{kbox}

  \textbf{PCP Arbitrage} (if parity violated):

  \textbf{$PV_A < PV_C$} ($C+dK < P+S_0$): Port.\ A cheap --- buy call, short put, short stock, invest $S_0+P-C$.\\
  {\tiny\begin{tabular}{@{}llll@{}}
    \toprule
                         & $t{=}0$          & $S_T{>}K$         & $S_T{\le}K$       \\
    \midrule
    Long call            & $-C$             & $+(S_T{-}K)$      & $0$               \\
    Short put            & $+P$             & $0$               & $-(K{-}S_T)$      \\
    Short stock          & $+S_0$           & $-S_T$            & $-S_T$            \\
    Invest $S_0{+}P{-}C$ & $-(S_0{+}P{-}C)$ & $+Ie^{rT}$        & $+Ie^{rT}$        \\
    \midrule
    \textbf{Total}       & $0$              & $Ie^{rT}{-}K{>}0$ & $Ie^{rT}{-}K{>}0$ \\
    \bottomrule
  \end{tabular}}\\
  where $I=S_0+P-C>dK$, so $Ie^{rT}>K$.

  \textbf{$PV_A > PV_C$} ($C+dK > P+S_0$): Port.\ C cheap --- short call, buy put, buy stock, borrow $S_0+P-C$.\\
  {\tiny\begin{tabular}{@{}llll@{}}
    \toprule
                         & $t{=}0$          & $S_T{>}K$         & $S_T{\le}K$       \\
    \midrule
    Short call           & $+C$             & $-(S_T{-}K)$      & $0$               \\
    Long put             & $-P$             & $0$               & $+(K{-}S_T)$      \\
    Long stock           & $-S_0$           & $+S_T$            & $+S_T$            \\
    Borrow $S_0{+}P{-}C$ & $+(S_0{+}P{-}C)$ & $-Ie^{rT}$        & $-Ie^{rT}$        \\
    \midrule
    \textbf{Total}       & $0$              & $K{-}Ie^{rT}{>}0$ & $K{-}Ie^{rT}{>}0$ \\
    \bottomrule
  \end{tabular}}\\
  where $I=S_0+P-C<dK$, so $Ie^{rT}<K$. Net profit $=K-Ie^{rT}>0$ in all cases.

  \subsection{Trading Strategies}
  \textbf{Bull spread:} Long $C_{K_1}$ + Short $C_{K_2}$, $K_1<K_2$. Capitalise on price increase, cap max profit to reduce cost \& risk.\\
  {\tiny\begin{tabular}{@{}llll@{}}
    \toprule
    $S_T$ range   & Long $C_{K_1}$ & Short $C_{K_2}$ & Total     \\
    \midrule
    $S_T\le K_1$  & $0$            & $0$             & $0$       \\
    $K_1<S_T<K_2$ & $S_T-K_1$      & $0$             & $S_T-K_1$ \\
    $S_T\ge K_2$  & $S_T-K_1$      & $K_2-S_T$       & $K_2-K_1$ \\
    \bottomrule
  \end{tabular}}

  \textbf{Bear spread:} Short $C_{K_1}$ + Long $C_{K_2}$ \emph{or} Long $P_{K_2}$ + Short $P_{K_1}$. Profit if $S_T\downarrow$.

  \textbf{Butterfly} ($K_1{<}K_2{<}K_3$, $2K_2{=}K_1{+}K_3$): Long $C_{K_1}$, Short $2C_{K_2}$, Long $C_{K_3}$. Low-vol bet ($S_T\approx K_2$).\\
  {\tiny\begin{tabular}{@{}lllll@{}}
    \toprule
    $S_T$ range      & $C_{K_1}$ & $C_{K_3}$ & $-2C_{K_2}$   & Total     \\
    \midrule
    $S_T\le K_1$     & $0$       & $0$       & $0$           & $0$       \\
    $K_1<S_T\le K_2$ & $S_T-K_1$ & $0$       & $0$           & $S_T-K_1$ \\
    $K_2<S_T< K_3$   & $S_T-K_1$ & $0$       & $-2(S_T-K_2)$ & $K_3-S_T$ \\
    $S_T\ge K_3$     & $S_T-K_1$ & $S_T-K_3$ & $-2(S_T-K_2)$ & $0$       \\
    \bottomrule
  \end{tabular}}

  \textbf{Straddle:} Long call + Long put (same $K$, $T$). Profit if $|S_T-K|$ large.

  \textbf{Covered call:} Long stock + Short call.\quad\textbf{Protective put:} Long stock + Long put.

  \subsection{Binomial Option Pricing}
  Single period: $S\to uS$ or $dS$, risk-free rate $r$, $R=1+r$. No-arb requires $u>R>d$.

  \textbf{Risk-neutral probability:} $q=\dfrac{R-d}{u-d}$, $1-q=\dfrac{u-R}{u-d}$

  \begin{kbox}
    \[C=\frac{1}{R}[q\,C_u+(1-q)\,C_d]\]
  \end{kbox}
  $C_u=\max(uS-K,0)$, $C_d=\max(dS-K,0)$ for call; $\max(K-\cdot,0)$ for put.

  \textbf{Multi-period:} backward induction. Each node same formula. \textbf{CRR:} $u=e^{\sigma\sqrt{\Delta t}}$, $d=1/u$.

  \subsection{Replication Portfolio (Binomial)}
  Construct portfolio: $x$ shares + $b$ in risk-free cash. Payoffs:
  \[\text{up: }ux+Rb=C_u\qquad\text{down: }dx+Rb=C_d\]
  Solve:
  \begin{kbox}
    \[x=\frac{C_u-C_d}{u-d}\qquad b=\frac{uC_d-dC_u}{R(u-d)}\]
    \[C=x+b=\frac{qC_u+(1-q)C_d}{R}\quad\checkmark\]
  \end{kbox}
  \textbf{$C < x+b$:} Short $x$ shares, borrow $b$, buy call. Net at $t{=}0$: $(x+b)-C>0$. At $T$: portfolio payoff $-C_u$ or $-C_d$ cancels call payoff, leftover $= R((x+b)-C)>0$.

  \textbf{$C > x+b$:} Short call, long $x$ shares, invest $b$. Net at $t{=}0$: $C-(x+b)>0$. At $T$: call payoff $-C_u$ or $-C_d$ cancels portfolio, leftover $= R(C-(x+b))>0$.


\end{multicols*}
\end{document}
